
\section{$A_\infty$ 双模与弗洛尔上同调}
\label{sec:3-7}

本节利用前面各节给出的构造，形变标准的弗洛尔“上边界”算子。我们先在第
3.7.1 小节准备滤过 $A_\infty$ 双模结构的代数框架，再在第 3.7.2 小节讨论
滤过 $A_\infty$ 双模同态。第 3.7.3 小节将在左右两个 $A_\infty$ 代数
（弱）无障碍时，代数地定义滤过 $A_\infty$ 双模上的形变上边界算子。

随后，第 3.7.4 小节对一对横截且相对 spin 的拉格朗日子流形
$(L^{(0)},L^{(1)})$ 作几何实现，并构造其弗洛尔上同调
\[
  HF\bigl((L^{(1)},b_1),(L^{(0)},b_0);\Lambda_{0,\mathrm{nov}}\bigr).
\]
第 3.7.5 小节处理 $L^{(0)}$ 与 $L^{(1)}$ 干净相交的情形。第 3.7.6 小节
给出若干简单的弗洛尔上同调计算，并讨论这些例子中的（弱）无障碍条件；第
3.7.7 小节则定义弗洛尔上同调上的乘法结构。

\subsection{代数框架}
\label{subsec:3-7-1}

本小节引入若干不涉及拉格朗日子流形的代数概念。令
$(C_0,\mathfrak m^0)$、$(C_1,\mathfrak m^1)$ 为
$\Lambda_{0,\mathrm{nov}}$ 上的两个滤过 $A_\infty$ 代数。设
$\bigoplus_{m\in\mathbb Z}D^m$ 是分次自由滤过
$\Lambda_{0,\mathrm{nov}}$-模，其滤过 $F^\lambda D^m$ 与 (3.2.12) 中的
定义类似，并满足 (3.2.12.1)--(3.2.12.5) 的相同条件。以 $D$ 表示
$\bigoplus_{m\in\mathbb Z}D^m$ 关于这个滤过的完备化。

考虑一族次数为 $+1$ 的运算
\[
  \mathfrak n_{k_1,k_0}\colon
  B_{k_1}(C_1[1])\otimes_{\Lambda_{0,\mathrm{nov}}}D[1]
  \otimes_{\Lambda_{0,\mathrm{nov}}}B_{k_0}(C_0[1])
  \longrightarrow D[1],
\]
它们满足
\begin{equation}
\begin{aligned}
&\mathfrak n_{k_1,k_0}\Bigl(
  F^{\lambda_1}(C_1[1])^{m_1}\otimes\cdots\otimes
  F^{\lambda_{k_1}}(C_1[1])^{m_{k_1}}
  \otimes F^{\lambda_0}(D[1])^{m_0}\\
&\hspace{7em}{}\otimes
  F^{\lambda'_1}(C_0[1])^{m'_1}\otimes\cdots\otimes
  F^{\lambda'_{k_0}}(C_0[1])^{m'_{k_0}}
  \Bigr)\\
&\qquad\subset
  F^{\sum_i\lambda_i+\lambda_0+\sum_i\lambda'_i}
  (D[1])^{\sum_i m_i+m_0+\sum_i m'_i+1}.
\end{aligned}
\tag{3.7.1}\label{eq:3-7-1}
\end{equation}

现在利用条件~\eqref{eq:3-7-1}，把 $\mathfrak n_{k_1,k_0}$ 延拓到
\[
  \widehat B(C_1[1])\widehat\otimes_{\Lambda_{0,\mathrm{nov}}}D[1]
  \widehat\otimes_{\Lambda_{0,\mathrm{nov}}}\widehat B(C_0[1]),
\]
即
\[
  \bigoplus_{k_1,k_0}B_{k_1}(C_1[1])
  \otimes_{\Lambda_{0,\mathrm{nov}}}D[1]
  \otimes_{\Lambda_{0,\mathrm{nov}}}B_{k_0}(C_0[1])
\]
的完备化。为此先注意，上述完备化是左
$\widehat B(C_1[1])$、右 $\widehat B(C_0[1])$ 的（形式）双余模。继而把
$\mathfrak n_{k_1,k_0}$ 延拓为双余导子\liftIndex{I124}
\[
\begin{aligned}
  \widehat d\colon{}&
  \widehat B(C_1[1])\widehat\otimes_{\Lambda_{0,\mathrm{nov}}}D[1]
  \widehat\otimes_{\Lambda_{0,\mathrm{nov}}}\widehat B(C_0[1])\\
  &\longrightarrow
  \widehat B(C_1[1])\widehat\otimes_{\Lambda_{0,\mathrm{nov}}}D[1]
  \widehat\otimes_{\Lambda_{0,\mathrm{nov}}}\widehat B(C_0[1]),
\end{aligned}
\]
其定义为
\begin{equation}
\begin{aligned}
&\widehat d(x_{1,1}\otimes\cdots\otimes x_{1,k_1}
  \otimes y\otimes x_{0,1}\otimes\cdots\otimes x_{0,k_0})\\
&={}
\sum_{\substack{k'_1\leq k_1\\k'_0\leq k_0}}
(-1)^{\deg x_{1,1}+\cdots+\deg x_{1,k_1-k'_1}+k_1-k'_1}
  x_{1,1}\otimes\cdots\otimes x_{1,k_1-k'_1}\\
&\qquad\otimes
  \mathfrak n_{k'_1,k'_0}
  (x_{1,k_1-k'_1+1},\ldots,x_{1,k_1},y,
   x_{0,1},\ldots,x_{0,k'_0})
  \otimes x_{0,k'_0+1}\otimes\cdots\otimes x_{0,k_0}\\
&\quad+
  \widehat d^{1}(x_{1,1}\otimes\cdots\otimes x_{1,k_1})
  \otimes y\otimes x_{0,1}\otimes\cdots\otimes x_{0,k_0}\\
&\quad+
  (-1)^{\sum_i\deg x_{1,i}+\deg y+k_1+1}
  x_{1,1}\otimes\cdots\otimes x_{1,k_1}\otimes y
  \otimes\widehat d^{0}(x_{0,1}\otimes\cdots\otimes x_{0,k_0}).
\end{aligned}
\tag{3.7.2}\label{eq:3-7-2}
\end{equation}
这里 $\widehat d^{i}$ 由 $\mathfrak m^i$ 定义（$i=0,1$）。采用约定
(Conv.2) (1.36) 时，可以把~\eqref{eq:3-7-2} 简写为
\begin{equation}
\begin{aligned}
\widehat d(\mathbf x_1\otimes y\otimes\mathbf x_0)
={}&\sum_{a,b}(-1)^{\deg'\mathbf x'_{1,a}}
  \mathbf x'_{1,a}\otimes
  \mathfrak n(\mathbf x''_{1,a}\otimes y\otimes\mathbf x'_{0,b})
  \otimes\mathbf x''_{0,b}\\
&+\widehat d^{1}\mathbf x_1\otimes y\otimes\mathbf x_0
 +(-1)^{\deg'\mathbf x_1+\deg' y}
  \mathbf x_1\otimes y\otimes\widehat d^{0}\mathbf x_0.
\end{aligned}
\tag{3.7.3}\label{eq:3-7-3}
\end{equation}
这里 $\deg'$ 是第 3.3 节所用的移位次数。条件~\eqref{eq:3-7-1} 保证
式~\eqref{eq:3-7-2} 延拓到上述完备化。

由于 $D$ 是自由 $\Lambda_{0,\mathrm{nov}}$-模，存在自由 $R$-模
$\overline D$ 以及同构
\[
  D\cong\overline D\otimes_R\Lambda_{0,\mathrm{nov}}.
\]
由
$D/\Lambda^+_{0,\mathrm{nov}}\cong\overline D[e,e^{-1}]$，运算
$\mathfrak n_{k_1,k_0}$ 诱导同态
\[
\begin{aligned}
\overline{\mathfrak n}_{k_1,k_0}\colon{}
&B_{k_1}(\overline C_1[1])\otimes_R\overline D[1]
 \otimes_RB_{k_0}(\overline C_0[1])\otimes_RR[e,e^{-1}]\\
&\longrightarrow\overline D[1]\otimes_RR[e,e^{-1}].
\end{aligned}
\]
我们假定 $\overline{\mathfrak n}_{k_1,k_0}$ 由一个运算
\[
  B_{k_1}(\overline C_1[1])\otimes_R\overline D[1]
  \otimes_RB_{k_0}(\overline C_0[1])
  \longrightarrow\overline D[1]
\]
诱导，并仍以 $\overline{\mathfrak n}_{k_1,k_0}$ 表示这个运算。

下面的条件对应于条件 3.2.18。

\begin{statement}{条件}{3.7.4}\label{cond:3-7-4}
给定元素 $\mathbf e_i\in C_i[1]$（$i=0,1$），有
\begin{equation}
\mathfrak n_{k_1,k_0}
(x_1\otimes\cdots\otimes\mathbf e_1\otimes\cdots\otimes x_{k_1-1}
 \otimes y\otimes x'_1\otimes\cdots\otimes x'_{k_0})=0
\tag{3.7.4.1}\label{eq:3-7-4-1}
\end{equation}
当 $k_1+k_0\geq2$，并且
\begin{equation}
\mathfrak n_{k_1,k_0}
(x_1\otimes\cdots\otimes x_{k_1}\otimes y
 \otimes x'_1\otimes\cdots\otimes\mathbf e_0\otimes\cdots
 \otimes x'_{k_0-1})=0
\tag{3.7.4.2}\label{eq:3-7-4-2}
\end{equation}
当 $k_1+k_0\geq2$。此外，
\begin{equation}
  \mathfrak n_{1,0}(\mathbf e_1,y)
  =(-1)^{\deg y}\mathfrak n_{0,1}(y,\mathbf e_0)=y.
\tag{3.7.4.3}\label{eq:3-7-4-3}
\end{equation}
\end{statement}

\begin{definition}{3.7.5}\label{def:3-7-5}
\begin{enumerate}[label=\textup{(\arabic*)},leftmargin=2.8em]
  \item 若 $\mathfrak n_{k_1,k_0}$ 满足~\eqref{eq:3-7-1} 和
  $\widehat d\circ\widehat d=0$，就称它在 $D$ 上定义了左
  $(C_1,\mathfrak m^1)$、右 $(C_0,\mathfrak m^0)$ 的滤过
  $A_\infty$ 双模结构；也简称 $D$ 是一个 $(C_1,C_0)$-滤过
  $A_\infty$ 双模\liftIndex{I115}\liftIndex{I116}，或称 $D$ 是
  $(C_1,C_0)$ 上的滤过 $A_\infty$ 双模。

  \item 假设 $(C_i,\mathfrak m^i)$ 具有幺元 $\mathbf e_i$。若条件 3.7.4
  成立，就称 $\{\mathbf e_1,\mathbf e_0\}$ 作为幺元起作用；此时有时也称
  $D$ 为\emph{含幺滤过 $A_\infty$ 双模}\liftIndex{I118}。

  \item 假设 $(C_i,\mathfrak m^i)$ 具有同伦幺元 $\mathbf e_i$。若 $D$
  上存在左 $(C_1^+,\mathfrak m^{1,+})$、右
  $(C_0^+,\mathfrak m^{0,+})$ 的滤过 $A_\infty$ 双模结构，使
  $\{\mathbf e_1^+,\mathbf e_0^+\}$ 作为幺元起作用，就称
  $\{\mathbf e_1,\mathbf e_0\}$ 作为同伦幺元起作用。这里如 (3.3.1)，
  \[
    C_i^+=C_i\oplus\Lambda_{0,\mathrm{nov}}\mathbf e_i^+
      \oplus\Lambda_{0,\mathrm{nov}}\mathbf f_i,
  \]
  并且由 (3.3.5.2)，
  \[
    \mathbf e_i^+=\mathfrak m_1^{i,+}(\mathbf f_i)+\mathbf e_i
      -\mathfrak h_2^i(1\mathbin{\overline\otimes}1).
  \]
  此时也称 $D$ 为\emph{同伦含幺 $A_\infty$ 双模}。\translatornote{原书把这条等式指向 (3.3.6)，但 (3.3.6) 定义的是
  $\overline{\mathfrak h}_k$；该等式由 (3.3.5.2) 直接移项得到。}

  \item 对 $R$ 上的无滤过 $A_\infty$ 代数
  $(\overline C_i,\overline{\mathfrak m}^{i})$（$i=0,1$），以类似方式
  定义 $R$ 上的 $(\overline C_1,\overline C_0)$-$A_\infty$ 双模
  $(\overline D,\overline{\mathfrak n})$；无滤过的含幺（同伦含幺）
  $A_\infty$ 双模也按上面 (2)、(3) 定义。

  若 $D$ 是 $(C_1,C_0)$-滤过 $A_\infty$ 双模，则 $\overline D$ 是
  $(\overline C_1,\overline C_0)$-$A_\infty$ 双模，称为 $D$ 的
  $R$-\emph{约化}\liftIndex{I117}。
  \translatornote{原书在这里把 $\overline D$ 称为 $A_\infty$ 单模；但它同时保留左右两个代数及全部
  $\overline{\mathfrak n}_{k_1,k_0}$ 运算，源索引也把本项列为 $A_\infty$ 双模的
  $R$-约化，故译文补回“bi-”。}

  \item 令 $G_i\subset\mathbb R_{\geq0}\times2\mathbb Z$ 为满足条件
  3.1.6 的幺半群，并令 $(C_i,\mathfrak m^i)$ 为 $G_i$-有隙滤过
  $A_\infty$ 代数。再令 $G\subset\mathbb R_{\geq0}\times2\mathbb Z$
  为满足条件 3.1.6、且包含 $G_1,G_0$ 的幺半群。若对每个
  $k_1,k_0$ 和 $\beta\in G$ 都存在
  \[
    \mathfrak n_{k_1,k_0,\beta}\colon
    B_{k_1}\overline C_1[1]\otimes\overline D[1]
    \otimes B_{k_0}\overline C_0[1]\longrightarrow\overline D[1],
  \]
  使得
  \[
    \mathfrak n_{k_1,k_0}
    =\sum_{\beta\in G}T^{\lambda(\beta)}e^{\mu(\beta)/2}
      \mathfrak n_{k_1,k_0,\beta},
  \]
  就称滤过 $A_\infty$ 双模 $(D,\mathfrak n_{k_1,k_0})$ 是
  $G$-\emph{有隙的}\liftIndex{I123}。
\end{enumerate}

与注 3.6.15 一样，若 $G'$ 是另一个满足条件 3.1.6 且包含上述 $G$ 的
幺半群，则每个 $G$-有隙滤过 $A_\infty$ 双模都是 $G'$-有隙的。
\translatornote{原书此处误引条件 3.7.4；该条件只刻画双模幺元，而本定义前文对
$G_i$ 和 $G$ 均要求幺半群满足条件 3.1.6，注 3.6.15 的对应论证亦使用后者。}
不必指明 $G$ 时，简称 $(D,\mathfrak n_{k_1,k_0})$ 为有隙滤过
$A_\infty$ 双模。
\end{definition}

\begin{example}{3.7.6}\label{ex:3-7-6}
令 $(C,\mathfrak m)$ 为滤过 $A_\infty$ 代数。可以把 $C$ 看作左、右均为
$(C,\mathfrak m)$ 的滤过 $A_\infty$ 双模：定义
\[
  \mathfrak n_{k_1,k_0}(\mathbf x\otimes y\otimes\mathbf z)
  =\mathfrak m_{k_1+k_0+1}(\mathbf x,y,\mathbf z),
\]
其中 $\mathbf x\in B_{k_1}C[1]$、$y\in C[1]$、
$\mathbf z\in B_{k_0}C[1]$。若 $C$ 有幺元（相应地，同伦幺元），它就在
$C$（相应地，$C^+$）上作为幺元（相应地，同伦幺元）起作用。若作为滤过
$A_\infty$ 代数的 $C$ 是有隙的，则作为滤过 $A_\infty$ 双模的 $C$ 也是
有隙的。
\end{example}

讨论滤过 $A_\infty$ 双模时，我们也将使用 $\Lambda_{\mathrm{nov}}$ 作为
系数环。这是因为，一对不同拉格朗日子流形的弗洛尔上同调在哈密顿同痕下只有
取 $\Lambda_{\mathrm{nov}}$ 系数时才保持同构（见第 5.3 节）。所谓取
$\Lambda_{\mathrm{nov}}$ 系数的滤过 $A_\infty$ 双模，是指一个定义在
$\Lambda_{0,\mathrm{nov}}$ 上的滤过 $A_\infty$ 双模与
$\Lambda_{\mathrm{nov}}$ 在 $\Lambda_{0,\mathrm{nov}}$ 上的张量积（见第
5.2.1 小节）。

\subsection{$A_\infty$ 双模同态}
\label{subsec:3-7-2}

令 $C_i,C'_i$ 为滤过 $A_\infty$ 代数（$i=0,1$），并令 $D,D'$ 分别为
$(C_1,C_0)$、$(C'_1,C'_0)$-滤过 $A_\infty$ 双模。设
$\mathfrak f^i\colon C_i\to C'_i$ 是滤过 $A_\infty$ 代数同态。

\begin{definition}{3.7.7}\label{def:3-7-7}
称一族 $\Lambda_{0,\mathrm{nov}}$-模同态
$\varphi=\{\varphi_{k_1,k_0}\}$ 是在同态对
$(\mathfrak f^1,\mathfrak f^0)$ 之上的\emph{滤过 $A_\infty$ 双模同态}
$D\to D'\liftIndex{I119}$，并简记为 $\varphi\colon D\to D'$，若
\[
\begin{aligned}
\varphi_{k_1,k_0}\colon{}
&B_{k_1}(C_1[1])\widehat\otimes_{\Lambda_{0,\mathrm{nov}}}D[1]
 \widehat\otimes_{\Lambda_{0,\mathrm{nov}}}B_{k_0}(C_0[1])\\
&\longrightarrow D'[1]
\end{aligned}
\]
满足以下三个性质：
\begin{enumerate}[label=\textup{(\arabic*)},leftmargin=2.8em]
  \item 它保持能量滤过，即
  \[
  \begin{aligned}
  \varphi_{k_1,k_0}\bigl(&F^{\lambda_1}B_{k_1}(C_1[1])
    \widehat\otimes F^\lambda D[1]
    \widehat\otimes F^{\lambda_0}B_{k_0}(C_0[1])\bigr)\\
  &\subset F^{\lambda_1+\lambda+\lambda_0}D'[1].
  \end{aligned}
  \]

  \item 令
  \[
  \begin{aligned}
  \widehat\varphi\colon{}
  &\widehat B(C_1[1])\widehat\otimes D[1]\widehat\otimes
    \widehat B(C_0[1])\\
  &\longrightarrow
    \widehat B(C'_1[1])\widehat\otimes D'[1]\widehat\otimes
    \widehat B(C'_0[1])
  \end{aligned}
  \]
  为 $\varphi_{k_1,k_0}$ 与 $\mathfrak f^i$ 诱导的余模同态。更明确地，采用
  约定 (Conv.2) (1.36)，对 $\mathbf x_i\in B(C_i[1])$、$y\in D[1]$，
  定义
  \begin{equation}
  \begin{aligned}
  \widehat\varphi(\mathbf x_1\otimes y\otimes\mathbf x_0)
  =\sum_{a,b}{}&\widehat{\mathfrak f}^{1}(\mathbf x'_{1,a})
  \otimes\varphi(\mathbf x''_{1,a}\otimes y\otimes\mathbf x'_{0,b})\\
  &\otimes\widehat{\mathfrak f}^{0}(\mathbf x''_{0,b}),
  \end{aligned}
  \tag{3.7.8}\label{eq:3-7-8}
  \end{equation}
  其中 $\widehat{\mathfrak f}^{i}$ 如 (3.2.28) 由 $\mathfrak f^i$ 诱导。
  还要求下列图式交换：

  \begin{center}
  \resizebox{\linewidth}{!}{$
  \begin{array}{ccc}
  \widehat B(C_1[1])\widehat\otimes D[1]\widehat\otimes\widehat B(C_0[1])
  &\xrightarrow{\ \widehat d\ }&
  \widehat B(C_1[1])\widehat\otimes D[1]\widehat\otimes\widehat B(C_0[1])
  \\
  {\scriptstyle\widehat\varphi}\big\downarrow
  &&\big\downarrow{\scriptstyle\widehat\varphi}
  \\
  \widehat B(C'_1[1])\widehat\otimes D'[1]\widehat\otimes\widehat B(C'_0[1])
  &\xrightarrow{\ \widehat d'\ }&
  \widehat B(C'_1[1])\widehat\otimes D'[1]\widehat\otimes\widehat B(C'_0[1]).
  \end{array}
  $}
  \par\smallskip\textbf{图式 3.7.1}
  \end{center}
  这里 $\widehat d$、$\widehat d'$ 依~\eqref{eq:3-7-2} 定义，所有完备张量积
  都取在 $\Lambda_{0,\mathrm{nov}}$ 上。

  \item 由 (1)，$\varphi_{k_1,k_0}$ 诱导
  \[
  \begin{aligned}
  \overline\varphi_{k_1,k_0}\colon{}
  &B_{k_1}(\overline C_1[1])\otimes_R\overline D[1]
   \otimes_RB_{k_0}(\overline C_0[1])\otimes_RR[e,e^{-1}]\\
  &\longrightarrow\overline D'[1]\otimes_RR[e,e^{-1}].
  \end{aligned}
  \]
  假定它由
  \[
    B_{k_1}(\overline C_1[1])\otimes_R\overline D[1]
    \otimes_RB_{k_0}(\overline C_0[1])\longrightarrow\overline D'[1]
  \]
  诱导，并仍以 $\overline\varphi_{k_1,k_0}$ 表示后一个映射。
\end{enumerate}

无滤过 $A_\infty$ 双模同态
$\overline\varphi\colon\overline D\to\overline D'$ 以类似方式定义，并要求
$\widehat{\overline\varphi}\circ\widehat{\overline d}
=\widehat{\overline d}'\circ\widehat{\overline\varphi}$。
\end{definition}

下面讨论有隙条件。

\begin{definition}{3.7.9}\label{def:3-7-9}
假设 $C_i,C'_i$ 是有隙滤过 $A_\infty$ 代数，且
$\mathfrak f^i\colon C_i\to C'_i$ 是 $G_i$-有隙滤过 $A_\infty$ 同态
（见定义 3.2.26、3.2.29）。令
$G\subset\mathbb R_{\geq0}\times2\mathbb Z$ 为满足条件 3.1.6 并包含
$G_1,G_0$ 的幺半群；令 $D$（相应地，$D'$）为左 $C_1$、右 $C_0$
（相应地，左 $C'_1$、右 $C'_0$）的 $G$-有隙滤过 $A_\infty$ 双模。
若存在
\[
  \varphi_{k_1,k_0,i}\colon
  B_{k_1}\overline C_1[1]\otimes\overline D[1]
  \otimes B_{k_0}\overline C_0[1]\longrightarrow\overline D'[1]
\]
使得
\[
  \varphi_{k_1,k_0}
  =\sum_iT^{\lambda_i}e^{\mu_i/2}\varphi_{k_1,k_0,i},
\]
就称同态对 $(\mathfrak f^1,\mathfrak f^0)$ 之上的滤过 $A_\infty$ 双模同态
$\varphi=\{\varphi_{k_1,k_0}\}\colon D\to D'$ 是
$G$-\emph{有隙的}\liftIndex{I120}。这里把 $D$ 识别为
$\overline D\otimes_R\Lambda_{0,\mathrm{nov}}$，其余对象亦然。

若 $G'$ 是另一个满足条件 3.1.6 且包含 $G$ 的幺半群，则每个 $G$-有隙滤过
$A_\infty$ 双模同态也是 $G'$-有隙的。不必指明 $G$ 时，简称
$\varphi$ 是有隙的。
\end{definition}

滤过 $A_\infty$ 双模同态的复合可以用显然的方式定义。

如上一小节末尾所述，在研究一对拉格朗日子流形的弗洛尔上同调时，为了研究其
不变性，我们还会使用 $\Lambda_{\mathrm{nov}}$ 作为系数环。在这种情形下，将引入
弱滤过 $A_\infty$ 双模同态的概念\liftIndex{I378}，以及相应的某种有隙条件；见第 5 章的第
5.2.1 小节。

\liftIndexRangeStart{I121}
\begin{definition}{3.7.10}\label{def:3-7-10}
假设 $C_i,C'_i$ 分别具有幺元 $\mathbf e_i,\mathbf e'_i$，并且
$\{\mathbf e_1,\mathbf e_0\}$ 在 $D$ 上作为幺元起作用，
$\{\mathbf e'_1,\mathbf e'_0\}$ 在 $D'$ 上作为幺元起作用。再设
$\mathfrak f^i\colon C_i\to C'_i$ 是定义 3.3.11 意义下的含幺
$A_\infty$ 同态，而 $\varphi\colon D\to D'$ 是同态对
$(\mathfrak f^1,\mathfrak f^0)$ 之上的滤过 $A_\infty$ 双模同态。若
\begin{equation}
  \varphi(\mathbf x_1\otimes\mathbf e_1\otimes\mathbf x_2,v,\mathbf y)
  =\varphi(\mathbf x,v,
    \mathbf y_1\otimes\mathbf e_0\otimes\mathbf y_2)=0,
\tag{3.7.11}\label{eq:3-7-11}
\end{equation}
就称 $\varphi$ \emph{保持幺元}，或称它为\emph{含幺 $A_\infty$ 双模同态}。
这里 $v\in D$，$\mathbf x,\mathbf x_1,\mathbf x_2\in BC_1[1]$，
$\mathbf y,\mathbf y_1,\mathbf y_2\in BC_0[1]$。式~\eqref{eq:3-7-11}
也包括 $\mathbf x_1,\mathbf x_2,\mathbf x,\mathbf y_1,\mathbf y_2$ 或
$\mathbf y$ 等于 $1$ 的情形。例如，
\[
\begin{gathered}
  \varphi(\mathbf e_1,v,\mathbf y)
  =\varphi(1,v,\mathbf y_1\otimes\mathbf e_0\otimes\mathbf y_2)=0,\\
  \varphi(\mathbf e_1,v,1)=\varphi(1,v,\mathbf e_0)=0.
\end{gathered}
\]
无滤过版本以同样方式定义。
\end{definition}

下面假定 $C_i,C'_i$ 分别具有同伦幺元 $\mathbf e_i,\mathbf e'_i$，而且它们
分别在 $D,D'$ 上作为同伦幺元起作用。再设
$\mathfrak f^i\colon C_i\to C'_i$ 是同伦含幺 $A_\infty$ 同态，且
$\mathfrak f^{i,+}\colon C_i^+\to C_i^{\prime+}$ 是定义 3.3.13 中保持幺元
$\mathbf e_i^+,\mathbf e_i^{\prime+}$ 的延拓。

\begin{definition}{3.7.12}\label{def:3-7-12}
在上述情形下，若同态对 $(\mathfrak f^1,\mathfrak f^0)$ 之上的滤过
$A_\infty$ 双模同态 $\varphi\colon D\to D'$ 可以延拓为同态对
$(\mathfrak f^{1,+},\mathfrak f^{0,+})$ 之上的滤过 $A_\infty$ 双模同态
$\varphi^+$，并且 $\varphi^+$ 按定义 3.7.10 保持幺元，就称 $\varphi$
\emph{保持同伦幺元}，或称它为\emph{同伦含幺 $A_\infty$ 双模同态}。
无滤过版本也以同样方式定义。
\end{definition}
\liftIndexRangeEnd{I121}

\subsection{弱无障碍性与形变}
\label{subsec:3-7-3}

在第 3.6 节中，对无障碍滤过 $A_\infty$ 代数 $C$ 以及每一对
$b_i\in\widehat{\mathcal M}(C)$，我们定义了上边界算子
\[
  \delta_{b_1,b_0}\colon C[1]\longrightarrow C[1]
\]
（定义 3.6.6）。当 $C$ 弱无障碍时，对给定的
$b\in\widehat{\mathcal M}_{\mathrm{weak}}(C)$ 也类似地定义了上边界算子
$\delta_{b,b}$（定义 3.6.33）。本小节讨论其 $A_\infty$ 双模版本。
\translatornote{原书在本段把 3.6.33 称为“引理”，又在本小节末把它称为“注”；
该编号的实际题头是“定义”，译文两处均据此改正。}

令 $D$ 为左 $(C_1,\mathfrak m^1)$、右 $(C_0,\mathfrak m^0)$ 的滤过
$A_\infty$ 双模，其双模结构由 $\mathfrak n_{k_1,k_0}$ 给出。先来形变这个
滤过 $A_\infty$ 双模结构。

\begin{statement}{定义--引理}{3.7.13}\label{def-lem:3-7-13}
\begin{enumerate}[label=\textup{(\arabic*)},leftmargin=2.8em]
  \item 对任意满足
  \[
    b_i\in C_i[1]^0,
    \qquad b_i\equiv0\pmod{\Lambda^+_{0,\mathrm{nov}}},
    \qquad i=0,1,
  \]
  的 $b_i$，定义一族算子
  \[
  \begin{aligned}
  &{}^{b_1}\mathfrak n_{k_1,k_0}^{b_0}
  (x_{1,1},\ldots,x_{1,k_1},y,x_{0,1},\ldots,x_{0,k_0})\\
  &\quad=
  \mathfrak n(e^{b_1}x_{1,1}e^{b_1}\cdots x_{1,k_1}e^{b_1},
    y,e^{b_0}x_{0,1}e^{b_0}\cdots x_{0,k_0}e^{b_0}),
  \end{aligned}
  \]
  其中
  $x_{1,1}\otimes\cdots\otimes x_{1,k_1}\in B_{k_1}C_1[1]$、
  $y\in D[1]$，且
  $x_{0,1}\otimes\cdots\otimes x_{0,k_0}\in B_{k_0}C_0[1]$。
  于是它在 $D$ 上定义左
  $(C_1,\mathfrak m^{1,b_1})$、右
  $(C_0,\mathfrak m^{0,b_0})$ 的滤过 $A_\infty$ 双模结构（形变后的
  $A_\infty$ 代数见定义 3.6.9）。把它记作
  ${}^{b_1}\mathfrak n^{b_0}$，称为 $\mathfrak n$ 由 $b_1,b_0$ 所作的
  \emph{形变 $A_\infty$ 双模结构}\liftIndex{I122}。依~\eqref{eq:3-7-3}，
  还把 ${}^{b_1}\mathfrak n^{b_0}$ 所诱导的算子记作
  ${}^{b_1}\widehat d^{b_0}$。

  \item 特别地，对 $y\in D[1]$ 置
  \[
  \begin{aligned}
    \delta_{b_1,b_0}(y)
    &={} ^{b_1}\mathfrak n^{b_0}(y)
      =\mathfrak n(e^{b_1}ye^{b_0})\\
    &=\sum_{k_1,k_0}
      \mathfrak n_{k_1,k_0}
      (b_1,\ldots,b_1,y,b_0,\ldots,b_0).
  \end{aligned}
  \]
\end{enumerate}
\end{statement}

${}^{b_1}\mathfrak n^{b_0}$ 定义滤过 $A_\infty$ 双模结构的证明与命题
3.6.10 的证明类似，故略去。第 5.2.2 小节将讨论滤过 $A_\infty$ 双模同态的
形变。

现在假定 $(C_1,\mathfrak m^1)$ 与 $(C_0,\mathfrak m^0)$ 无障碍，并令
$b_1,b_0$ 分别为二者的定界余链。用 $b_1,b_0$ 形变 $A_\infty$ 双模结构，
得到如下上边界算子
\[
  \delta_{b_1,b_0}\colon D\longrightarrow D.
\]
在几何语境中，这对应于标准弗洛尔上边界算子的形变。

\begin{lemma}{3.7.14}\label{lem:3-7-14}
对 $(C_1,\mathfrak m^1)$、$(C_0,\mathfrak m^0)$ 的任意定界余链
$b_1,b_0$，都有
\[
  \delta_{b_1,b_0}\circ\delta_{b_1,b_0}=0.
\]
\end{lemma}

\begin{proof}
证明与引理 3.6.8 相同，故略。
\end{proof}

接着假定 $C_1,C_0$ 分别具有幺元 $\mathbf e_1,\mathbf e_0$，且
$\{\mathbf e_1,\mathbf e_0\}$ 作为幺元起作用。仍可像定义--引理 3.7.13
那样用 $b_1,b_0$ 形变滤过 $A_\infty$ 双模结构；容易看出，
$\{\mathbf e_1,\mathbf e_0\}$ 在形变后的滤过 $A_\infty$ 双模上仍作为幺元
起作用。

若 $C_i$ 弱无障碍，则对 $i=0,1$ 有定义 3.6.33 (2) 所定义的势函数
\begin{equation}
  \mathfrak{PO}_i\colon
  \widehat{\mathcal M}_{\mathrm{weak}}(C_i)
  \longrightarrow\Lambda_{0,\mathrm{nov}}^{+(0)}.
\tag{3.7.15}\label{eq:3-7-15}
\end{equation}
对每个 $b_i\in\widehat{\mathcal M}_{\mathrm{weak}}(C_i)$，定义
$\delta_{b_1,b_0}\colon D\to D$ 为
\begin{equation}
  \delta_{b_1,b_0}(y)=\mathfrak n(e^{b_1}ye^{b_0}).
\tag{3.7.16}\label{eq:3-7-16}
\end{equation}

\begin{proposition}{3.7.17}\label{prop:3-7-17}
对 $y\in D$，有
\[
  \delta_{b_1,b_0}\circ\delta_{b_1,b_0}(y)
  =\bigl(-\mathfrak{PO}_1(b_1)+\mathfrak{PO}_0(b_0)\bigr)ey.
\]
特别地，若
$\mathfrak{PO}_1(b_1)=\mathfrak{PO}_0(b_0)$，就可以定义
$(D,\delta_{b_1,b_0})$ 的上同调。
\end{proposition}

\begin{proof}
直接计算得到
\[
\begin{aligned}
\delta_{b_1,b_0}(\delta_{b_1,b_0}(y))
&=\mathfrak n(e^{b_1}\mathfrak n(e^{b_1}ye^{b_0})e^{b_0})\\
&=\mathfrak n\bigl(\widehat d(e^{b_1}ye^{b_0})\bigr)
  -\mathfrak n\bigl(\widehat d(e^{b_1})ye^{b_0}\bigr)
  +(-1)^{\deg y}\mathfrak n\bigl(e^{b_1}y\widehat d(e^{b_0})\bigr)\\
&=\bigl(-\mathfrak{PO}_1(b_1)+\mathfrak{PO}_0(b_0)\bigr)ey.
\end{aligned}
\]
最后一个等号使用了 (3.7.4.3)。
\end{proof}

\begin{remark}{3.7.18}\label{rem:3-7-18}
若 $C_i$ 具有作为同伦幺元起作用的同伦幺元 $\mathbf e_i$，并且按定义
3.6.29 (2) 的意义弱无障碍，则考虑如 (3.3.1) 的
\[
  C_i^+=C_i\oplus\Lambda_{0,\mathrm{nov}}\mathbf e_i^+
  \oplus\Lambda_{0,\mathrm{nov}}\mathbf f_i
\]
后，上述论证仍然成立。（注意，这种情形下
$\delta_{b_1,b_0}\colon D\to D$ 有定义；比较定义 3.6.33 (1)。）第
3.7.6 小节将给出命题 3.7.17 的一个几何例子。
\end{remark}
